\section{Introduction}
\label{sec:introduction}

Automated unit test generation using Large Language Models (LLMs) has emerged as a promising alternative to traditional Search-Based Software Testing (SBST) tools like EvoSuite~\cite{fraser2011evosuite}, producing human-readable JUnit tests with minimal boilerplate~\cite{schafer2024empirical, alshahwan2024automated}. However, while modern LLMs excel at generating syntactically compilable code, they frequently stumble upon \textbf{the test oracle problem}: test suites compile successfully but fail during execution due to erroneous assertions or unhandled runtime exceptions~\cite{yuan2024evaluating, gu2024testart, konstantinou2025yate}. Specifically, models frequently hallucinate expected outputs by guessing from method names or textbook algorithmic conventions rather than tracing the actual control flow and state mutations of the System Under Test (SUT).

This reliability bottleneck is exacerbated in \emph{class-level} object-oriented software. Unlike standalone, stateless functions evaluated in early benchmarks such as HumanEval~\cite{chen2021evaluating, siddiq2024using}, real-world OOP classes encapsulate mutable instance fields, constructor preconditions, and inter-method dependencies~\cite{du2024classeval, xue2025classevalt}. When an LLM fails to track sequential state transitions, it inevitably derives false oracles. While prior repair approaches predominantly rely on iterative compiler feedback to resolve syntax errors~\cite{yuan2024evaluating, gu2024testart, singh2026qualitative}, little empirical work has investigated whether structured, source-first \emph{prompt refinement} can resolve semantic oracle failures on class-level Java benchmarks without peeking at runtime logs or reference tests.

To address this challenge, we present an empirical investigation of failure patterns and propose \textbf{Oracle-Aware Prompt Refinement (OAPR)} for LLM-generated JUnit tests on the \textbf{ClassEval-T} Java benchmark~\cite{xue2025classevalt}. We evaluate two frontier models---ChatGPT and Gemini---across 94 class-level tasks spanning seven application domains under a standardized baseline prompt ($P_0$) and an OAPR refined prompt ($P_1$). Specifically, our investigation explores two central research questions: first, examining the baseline prevalence and underlying semantic failure taxonomy of LLM-generated test failures in stateful object-oriented software; and second, evaluating the extent to which structured, source-first prompt refinement can guide models to independently recover broken test oracles and identify the remaining failure boundaries.

This paper makes the following key contributions:
\begin{enumerate}
    \item \textbf{Controlled Class-Level Evaluation:} We benchmark ChatGPT and Gemini on 94 ClassEval-T classes (187 valid pairs after infrastructure quarantine), identifying 18 valid $P_0$ failures (11 ChatGPT, 7 Gemini) with baseline Combined Success rates of 88.30\% and 92.47\%.
    \item \textbf{Fine-Grained Oracle Failure Taxonomy:} We classify all 18 failures across seven domains into state-transition, transformation, boundary, and parsing oracle errors.
    \item \textbf{Oracle-Aware Prompt Refinement (OAPR):} We design $P_1$, a nine-step semantic prompt protocol that enforces strict precondition analysis, state tracking, and arithmetic recomputation.
    \item \textbf{Empirical Recovery Validation:} In a focused recovery setting, OAPR recovers \textbf{15/18 (83.33\%)} failed pairs (10/11 for ChatGPT, 5/7 for Gemini), elevating final success rates to 98.94\% and 97.85\%, while identifying the boundary limitations of prompt-only repair on complex spatial and multi-rail cipher tasks.
\end{enumerate}

The rest of this paper is organized as follows: Section~\ref{sec:background} surveys related work; Section~\ref{sec:dataset_setup} describes ClassEval-T and the experimental setup; Section~\ref{sec:methodology} outlines our methodology; Section~\ref{sec:results} presents the results; Section~\ref{sec:discussion} and Section~\ref{sec:threats} discuss implications and threats; and Section~\ref{sec:conclusion} concludes.
